\section{The Digital Negative: Characteristic Curve Model}
\label{sec:negative}

\subsection{Requirements on the Curve Model}

The transfer from log exposure to optical density is the mathematical heart of
the application. The model must satisfy six requirements:

\begin{enumerate}
\item \textbf{R1.} Reproduce the toe, straight-line, and shoulder regions of a
measured $D$--$\logE$ curve to within AC-1 tolerance.
\item \textbf{R2.} Provide \emph{independent} control of each region. A change
to shoulder sharpness must not move the speed point; a change to gamma must not
change $\Dmax$.
\item \textbf{R3.} Be $C^\infty$. Any discontinuity in the first derivative
produces a visible contour in a smooth gradient, and any discontinuity in the
second derivative is visible in a sky.
\item \textbf{R4.} Have a closed-form derivative, because the derivative
\emph{is} the local gamma, which the grain model
(Section~\ref{sec:grain}) requires pointwise.
\item \textbf{R5.} Be monotonically non-decreasing for all reachable
parameters (AC-6).
\item \textbf{R6.} Evaluate in a handful of GPU instructions.
\end{enumerate}

Spline-fit tables satisfy R1 and R6 but violate R2 catastrophically --- moving
one control point changes the curve globally in ways that do not correspond to
any physical parameter, which destroys parameter honesty. The classical
Hurter--Driffield treatment \cite{hurter1890,james1977} gives the straight line
but no closed form for the extremes. Generalized logistic (Richards) curves are $C^\infty$ and monotone but
couple toe and shoulder sharpness through a single asymmetry exponent, failing
R2.

\subsection{The Softplus Density Model}

\EM{} uses a formulation built from softplus primitives. Define the softplus
with sharpness parameter $a > 0$:

\begin{equation}
\softp{a}{u} \;=\; a \ln\!\left(1 + e^{u/a}\right).
\label{eq:softplus}
\end{equation}

This function is the canonical smooth relaxation of $\max(u, 0)$: it is
$C^\infty$, strictly increasing, satisfies $\softp{a}{u} \to \max(u,0)$ as $a
\to 0^+$, and has the logistic function as its derivative,
$\tfrac{d}{du}\softp{a}{u} = \sigma(u/a)$ with $\sigma(t) = (1+e^{-t})^{-1}$.

Let $x = \logE$ be the log exposure reaching a given layer, and define the
straight-line coordinate

\begin{equation}
u(x) \;=\; \gamma\,\big(x - x_{0}\big),
\label{eq:ucoord}
\end{equation}

where $\gamma$ is the straight-line gamma and $x_0$ is the \emph{extrapolated
speed point}: the log exposure at which the straight-line portion, extended
downward, intersects $\Dmin$.

The characteristic curve is then

\begin{equation}
\boxed{\;D(x) \;=\; \Dmin \;+\; \softp{\kappa_t}{u(x)} \;-\; \softp{\kappa_s}{u(x) - \Delta D}\;}
\label{eq:charcurve}
\end{equation}

with $\Delta D = \Dmax - \Dmin$ the density range, $\kappa_t > 0$ the toe
softness, and $\kappa_s > 0$ the shoulder softness.

\subsection{Behavior and Verification of the Model}

The three asymptotic regimes follow directly from \eqref{eq:softplus}:

\paragraph{Toe, $u \ll 0$} Both softplus terms vanish exponentially and $D \to
\Dmin$. Approach is exponential with scale $\kappa_t/\gamma$ in log exposure,
matching the observed sub-threshold behavior of real emulsions.

\paragraph{Straight line, $\kappa_t \ll u \ll \Delta D - \kappa_s$} The first
term is $\approx u$, the second $\approx 0$, giving

\begin{equation}
D(x) \approx \Dmin + \gamma\,(x - x_0),
\end{equation}

a line of slope exactly $\gamma$ intercepting $\Dmin$ at $x_0$, as required by
R2.

\paragraph{Shoulder, $u \gg \Delta D$} The first term is $\approx u$, the
second $\approx u - \Delta D$, and $D \to \Dmin + \Delta D = \Dmax$.

\medskip
The derivative --- the \emph{point gamma} $\gamma(x)$ --- is available in closed
form, satisfying R4:

\begin{equation}
\boxed{\;\gamma(x) \;=\; \frac{dD}{dx} \;=\; \gamma\left[\sigma\!\left(\frac{u}{\kappa_t}\right) - \sigma\!\left(\frac{u - \Delta D}{\kappa_s}\right)\right]\;}
\label{eq:pointgamma}
\end{equation}

Monotonicity (R5) follows immediately: $\sigma$ is strictly increasing, and
since $u > u - \Delta D$ whenever $\Delta D > 0$, we have $\sigma(u/\kappa_t) >
\sigma((u-\Delta D)/\kappa_s)$ provided $\kappa_t \le \kappa_s$ or, more
generally, provided $\Delta D$ is large relative to the softness parameters.
The precise condition is stated below.

\subsection{Well-Formedness Condition}

The two softplus terms overlap when $\Delta D$ is not large compared to
$\kappa_t + \kappa_s$, in which case the curve never reaches $\Dmax$ and, in the
extreme, can become non-monotone. Evaluating \eqref{eq:pointgamma} at the
midpoint $u = \Delta D/2$ gives the minimum-margin case; requiring
$\gamma(x) > 0$ everywhere yields the sufficient condition

\begin{equation}
\Delta D \;>\; \left(\kappa_t + \kappa_s\right)\ln\!\left(\frac{\kappa_s}{\kappa_t}\right)
\quad\text{when } \kappa_s > \kappa_t,
\label{eq:wellformed1}
\end{equation}

and monotonicity is unconditional when $\kappa_s \le \kappa_t$. In practice
\EM{} enforces the stricter engineering constraint

\begin{equation}
\Delta D \;\geq\; 4\left(\kappa_t + \kappa_s\right),
\label{eq:wellformed}
\end{equation}

which additionally guarantees that the attained maximum density is within
$0.02$ density units of the nominal $\Dmax$ and that a genuine straight-line
region of at least $\Delta D / 2$ exists. The parameter validator in
\texttt{EmulsionCore} rejects any profile violating \eqref{eq:wellformed} at
load time, and the UI clamps $\kappa$ sliders so that the condition cannot be
violated interactively.

\subsection{Per-Record Parameterization and the Orange Mask}

Each color record has an independent parameter set
$\theta_c = (\Dmin{}_{,c}, \Delta D_c, \gamma_c, x_{0,c}, \kappa_{t,c},
\kappa_{s,c})$, $c \in \{R,G,B\}$, so \eqref{eq:charcurve} is evaluated three
times.

The $\Dmin$ vector deserves specific attention because it encodes the
\emph{orange mask} of color negative film. Colored couplers in the magenta and
cyan layers absorb blue and green light even where undeveloped, producing the
familiar orange base. In Status M density \cite{iso5} this appears as

\begin{equation}
\Dmin^{\mathrm{(colorneg)}} \approx (0.58,\; 0.92,\; 1.28)^\top
\end{equation}

for a typical modern stock. The mask is not decorative: it is an
inter-image correction that compensates for unwanted absorptions of the cyan
and magenta dyes, and it is calibrated to be removed by the printing stage.
Modelling the mask but failing to model its removal produces the ``scanned
negative that was never printed'' look --- muddy, low-saturation, warm ---
which is a real and recognizable failure mode. Section~\ref{sec:print}
specifies the removal.

Because the mask density is largely independent of exposure, it can be carried
as a constant offset. \EM{} nevertheless models a small exposure dependence,
since real masking couplers are consumed in proportion to development:

\begin{equation}
\Dmin{}_{,c}(x) = \Dmin{}_{,c}^{(0)} - m_c \cdot \frac{D_c(x) - \Dmin{}_{,c}^{(0)}}{\Delta D_c},
\label{eq:maskdepletion}
\end{equation}

with $m_c$ the mask depletion coefficient, typically $(0.00, 0.06, 0.10)$.
Equation \eqref{eq:maskdepletion} is applied as a single fixed-point iteration
rather than solved exactly; the correction is small and one iteration is
below the AC-1 tolerance.

\subsection{Curve Family Illustration}

Figure~\ref{fig:charcurve} plots \eqref{eq:charcurve} for the three records of a
representative daylight color negative stock, and
Figure~\ref{fig:pointgamma} plots the corresponding point gamma
\eqref{eq:pointgamma}.

\begin{figure}[htbp]
\centering
\begin{tikzpicture}
\begin{axis}[
  width=\columnwidth, height=5.4cm,
  xlabel={$\logE$ (relative)},
  ylabel={Status M density $D$},
  xmin=-3, xmax=2, ymin=0, ymax=3.2,
  grid=both, grid style={draw=emLine!20},
  legend style={font=\scriptsize, at={(0.02,0.98)}, anchor=north west, draw=emLine!40},
  label style={font=\scriptsize}, tick label style={font=\scriptsize},
  samples=250,
]
% Blue record (highest Dmin, mask)
\addplot[emBlue, very thick, domain=-3:2]
  {1.28 + 0.13*ln(1+exp((0.66*(x+2.05))/0.13)) - 0.11*ln(1+exp((0.66*(x+2.05)-1.85)/0.11))};
\addlegendentry{Blue record}
% Green record
\addplot[emGreen, very thick, domain=-3:2]
  {0.92 + 0.14*ln(1+exp((0.63*(x+2.10))/0.14)) - 0.11*ln(1+exp((0.63*(x+2.10)-1.90)/0.11))};
\addlegendentry{Green record}
% Red record
\addplot[emAcc, very thick, domain=-3:2]
  {0.58 + 0.15*ln(1+exp((0.61*(x+2.15))/0.15)) - 0.12*ln(1+exp((0.61*(x+2.15)-1.95)/0.12))};
\addlegendentry{Red record}
\end{axis}
\end{tikzpicture}
\caption{Characteristic curves from \eqref{eq:charcurve} for a daylight color negative stock. The vertical separation at the base is the orange mask. Note that the three records are near-parallel through the straight-line region --- non-parallelism is \emph{crossover}, and it is the dominant cause of colors that cannot be corrected by a global balance change.}
\label{fig:charcurve}
\end{figure}

\begin{figure}[htbp]
\centering
\begin{tikzpicture}
\begin{axis}[
  width=\columnwidth, height=4.4cm,
  xlabel={$\logE$ (relative)},
  ylabel={point gamma $\gamma(x)$},
  xmin=-3, xmax=2, ymin=0, ymax=0.8,
  grid=both, grid style={draw=emLine!20},
  label style={font=\scriptsize}, tick label style={font=\scriptsize},
  samples=250,
]
\addplot[emAcc, very thick, domain=-3:2]
  {0.61*( 1/(1+exp(-(0.61*(x+2.15))/0.15)) - 1/(1+exp(-(0.61*(x+2.15)-1.95)/0.12)) )};
\addplot[emGreen, very thick, domain=-3:2]
  {0.63*( 1/(1+exp(-(0.63*(x+2.10))/0.14)) - 1/(1+exp(-(0.63*(x+2.10)-1.90)/0.11)) )};
\end{axis}
\end{tikzpicture}
\caption{Point gamma from \eqref{eq:pointgamma} for the red and green records. The plateau is the straight line; the rise and fall are the toe and shoulder. This function is consumed directly by the grain variance model \eqref{eq:grainvar} and by the interlayer inhibition weighting.}
\label{fig:pointgamma}
\end{figure}

\subsection{Derived Sensitometric Quantities}

The following quantities are computed from $\theta_c$ on the host and exposed
in the UI as read-only diagnostics. They are the standard vocabulary of
sensitometry and their presence is a substantial part of the product's claim to
professional seriousness.

\paragraph{Speed point} The log exposure producing density $\Dmin + 0.10$:

\begin{equation}
x_{\mathrm{sp}} = x_0 + \frac{\kappa_t}{\gamma}\ln\!\left(e^{0.10/\kappa_t} - 1\right).
\label{eq:speedpoint}
\end{equation}

This inverts the toe branch of \eqref{eq:charcurve} exactly, since the shoulder
term is negligible at this density. ISO speed follows as $S = 0.8 /
10^{x_{\mathrm{sp}}}$ in the arithmetic scale \cite{iso5800,iso6}.

\paragraph{Contrast index} Following standard sensitometric practice
\cite{kodakh740}, the slope of the chord
from the speed point to a point $2.00$ log-exposure units higher:

\begin{equation}
\mathrm{CI} = \frac{D(x_{\mathrm{sp}} + 2.00) - D(x_{\mathrm{sp}})}{2.00}.
\label{eq:ci}
\end{equation}

For a normal-process color negative $\mathrm{CI} \approx 0.55$--$0.62$; for
black-and-white pictorial negatives $0.50$--$0.58$.

\paragraph{Useful exposure latitude} The interval over which point gamma
exceeds a fraction $\lambda$ of its maximum (default $\lambda = 0.5$),
obtained by solving $\gamma(x) = \lambda\gamma_{\max}$ numerically with two
Newton steps from the analytic straight-line bracket. Reported in stops as
$\mathrm{L} = (x_{\mathrm{hi}} - x_{\mathrm{lo}})/\log_{10} 2$. Modern color
negative stocks report 12--14 stops by this measure, which is why the format
tolerates overexposure so gracefully --- and why the model, if correct,
reproduces that tolerance without a special case.

\subsection{Monochrome Stocks}

Monochrome emulsions use the same formulation with a single parameter set and a
spectral weighting applied before \eqref{eq:charcurve} to convert the
three-channel working image into the single exposure the panchromatic layer
receives:

\begin{equation}
E_{\mathrm{pan}} = \mathbf{w}^\top \mathbf{E}, \qquad
\mathbf{w} = \mathbf{w}_{\mathrm{spectral}} \odot \mathbf{w}_{\mathrm{filter}},
\label{eq:pan}
\end{equation}

where $\mathbf{w}_{\mathrm{spectral}}$ is the stock's normalized spectral
sensitivity integrated against the AP1 primaries and
$\mathbf{w}_{\mathrm{filter}}$ models an optional contrast filter (yellow,
orange, red, green) \cite{kodakh1}. This is the mechanism by which a red filter darkens a blue
sky in \EM{} for the same reason it does on a real camera, rather than by a
scripted sky-darkening effect. Default weights for a panchromatic stock are
$\mathbf{w}_{\mathrm{spectral}} = (0.30, 0.59, 0.11)$; an orthochromatic
profile uses $(0.02, 0.68, 0.30)$, which correctly renders red lips dark.

\subsection{Shader Implementation}

Equation~\eqref{eq:charcurve} costs two \texttt{log}, two \texttt{exp}, and
approximately ten fused multiply-adds per channel. The naive evaluation
overflows for large $u/\kappa$; the numerically stable form used in the shader
is

\begin{equation}
\softp{a}{u} = \max(u, 0) + a\ln\!\left(1 + e^{-|u|/a}\right),
\label{eq:softplus-stable}
\end{equation}

which is exact and never evaluates $e^{t}$ for $t > 0$.

\begin{lstlisting}[language=Metal, caption={Stable characteristic curve evaluation.}, label={lst:charcurve}]
inline float softplus_stable(float u, float a) {
    // sp_a(u) = max(u,0) + a*log1p(exp(-|u|/a))
    return max(u, 0.0f) + a * log1p(exp(-abs(u) / a));
}

inline float3 characteristic_curve(float3 logE, constant StockParams &p) {
    float3 u = p.gamma * (logE - p.x0);
    float3 sp_t, sp_s;
    for (uint c = 0; c < 3; ++c) {
        sp_t[c] = softplus_stable(u[c],              p.kappa_t[c]);
        sp_s[c] = softplus_stable(u[c] - p.dD[c],    p.kappa_s[c]);
    }
    return p.dMin + sp_t - sp_s;
}

inline float3 point_gamma(float3 logE, constant StockParams &p) {
    float3 u = p.gamma * (logE - p.x0);
    float3 a = 1.0f / (1.0f + exp(-u / p.kappa_t));
    float3 b = 1.0f / (1.0f + exp(-(u - p.dD) / p.kappa_s));
    return p.gamma * (a - b);
}
\end{lstlisting}

Note that \texttt{log1p} is required rather than \texttt{log(1+x)}: for large
$|u|/a$ the argument $e^{-|u|/a}$ underflows toward zero and $\log(1+x)$ loses
all significant digits, producing a stair-stepped toe that is plainly visible in
deep shadow. This is the kind of defect that survives casual review and is
caught only by the ramp test in NFR-06.
